% !PRINTED 480
\weberpage{480}
Auf dieselbe Weise sind die in der Tabelle am Ende aufgeführten
Klasseninvarianten für
\[
m=75,\quad 36,\quad 100,\quad 63,\quad 175
\]
berechnet, und diese Rechnungen lassen sich auch noch weiter
fortsetzen.

II. Aus den \gesperrt{Schlaeflischen} Modulargleichungen läßt sich noch auf
verschiedene andere Arten für die Berechnung von Klasseninvarianten
Nutzen ziehen.

Setzen wir
\[
\omega=-\frac{m}{\omega},\quad \omega=\sqrt{-m},
\]
so wird
\weq{11}{f(\omega)=f\Bigl(\frac{\omega}{m}\Bigr)=f(\sqrt{-m});}
wenn also in der Modulargleichung für den Transformationsgrad $m$
\[
u=v
\]
gesetzt wird, so ergibt sich eine Gleichung, unter deren Wurzeln
$u=f(\sqrt{-m})$ vorkommt. In dem System I. des § 73 ist dann
immer $A=2$ zu setzen, und es ergeben sich die Formeln:
\[
m=3,\quad B=u^6-\frac{8}{u^6}=2,
\]
\[
m=5,\quad B=u^4-\frac{4}{u^4}=2,
\]
\[
m=7,\quad B=u^6+\frac{8}{u^6}=9,
\]
\[
m=11,\quad B=u^2-\frac{2}{u^2},\quad B^5-B^3-2B=2,
\]
\[
m=13,\quad B=u^{12}-\frac{64}{u^{12}}=9\cdot2^5,
\]
\[
m=17,\quad B=u^8+\frac{16}{u^8},\quad B^2-68B-544=0,
\]
\[
m=19,\quad B=u^6-\frac{8}{u^6},\quad B^3-38B^2+252B-648=0.
\]
Die Fälle $m=3,7,11,19$ ergeben keine neuen Resultate, können
aber zur Verifizierung der früher gefundenen verwandt werden;
aus $m=5,13$ erhalten wir durch Auflösung einer
quadratischen Gleichung:
%CONT>>
